\section{Other decompositions and the concluding examples}
\subsection{Binary triples and ternary triples}
For binary $J(n,3)$ with odd $n$, the simplex identity gives, for
$x=\{a,b,c\}$,
\[
 e_x=r_x^2+w_{\{a,b\}}+w_{\{a,c\}}+w_{\{b,c\}}.
\]
The first term is in $C=W_2^\perp$ and the remaining term is in $W_2$.
This is the first decomposition of source Section~6; summing against
the coordinates of any vector gives its orthogonal splitting. The
source also records a typographical correction in its predecessor [21];
that is historical attribution, not a newly checked claim about that
predecessor's printed pages.

For ternary triples and $n\equiv1\pmod3$, the second displayed identity is
\begin{equation}\label{eq:ternary-split}
 e_{\{a,b,c\}}=\mathbf j-r_x^0+r_x^1+w_a+w_b+w_c.
\end{equation}
At a triple meeting $x$ in $t=0,1,2,3$ points, the right side has
coefficient $1-[t=0]+[t=1]+t$, namely $0,0,0,1$ modulo three.
This proves the literal identity directly. The membership
$w_a\in C_0^\perp$ also follows by counting: for a triple disjoint
from $a$, the relevant count is $\binom{n-4}{2}=0$ modulo three,
and for a triple containing $a$ it is zero. The additional assertions
$r_x^1\in C_0$ and $\mathbf j\in C_0$ or $C_0^\perp$ are retained
from the source's explicit invocation of Key--Moori--Rodrigues [23],
Proposition~18. With those memberships, summing
\eqref{eq:ternary-split} gives exactly the source's splitting of every
vector. In particular $C_0+C_0^\perp$ is the whole space, so the
dimension argument proves LCD and makes that splitting orthogonal and
unique. This application does not present that borrowed membership
theorem as a consequence of the unrelated $n\equiv5\pmod9$ case.

\subsection{A uniform odd-graph correction}
\begin{proposition}\label{prop:odd-graph}
For every $k\geq1$, let $A$ be the binary adjacency matrix of
$KG(2k+1,k)$ and put $E=u_{k-1}d_k$. Then
\[
 A^2=I+E,\qquad A^3=A,\qquad
 C_A=W_{k-1}^\perp,\qquad C_A^\perp=W_{k-1}.
\]
In particular $C_A$ is LCD. Its orthogonal decomposition on a unit
vector is
\begin{equation}\label{eq:odd-split}
 e_x=\sum_{y\cap x=\varnothing,\ |y|=k}r_y
          +\sum_{R\subset x,\ |R|=k-1}w_R.
\end{equation}
Moreover $\mathbf j\in C_A$ for even $k$ and
$\mathbf j\in C_A^\perp$ for odd $k$.
\end{proposition}
\begin{proof}
For distinct $k$-sets with intersection size $j$, the common-neighbour
count is $\binom{1+j}{k}$: it is one precisely at $j=k-1$.
The diagonal count is $k+1$. Thus
$A^2=J(2k+1,k)_{\rm adj}+(k+1)I=I+E$ by the common Johnson identity.
For a fixed $(k-1)$-set $R$, the $T$-coordinate of $Aw_R$ counts
$k$-sets containing $R$ and disjoint from $T$. If $R\cap T\ne\varnothing$
it is zero; otherwise there are exactly two choices for the remaining
point. Therefore $Au_{k-1}=0$, and $AE=0$. Multiplying the first
identity by $A$ proves $A^3=A$.

It follows that $\im A=\im A^2$: one inclusion is automatic and
$Av=A^2(Av)$ proves the other. Since $2k+1$ is odd,
\Uref{ULCD.Johnson} makes $E$ the orthogonal projector onto
$W_{k-1}$ and $I+E$ the projector onto its dual. This proves the
code identities and LCD. Applying $A^2+E=I$ to $e_x$ gives
\eqref{eq:odd-split}; its two summands lie in the asserted orthogonal
spaces. Finally $d_k\mathbf j=(k+2)\mathbf j$ and
$u_{k-1}\mathbf j=k\mathbf j$, so $E\mathbf j=k\mathbf j$ in
characteristic two. This proves the constant-word clause.
\end{proof}
This also gives dimension $\binom{2k}{k}$ and minimum distance $k+1$,
with all minimum words the $(k+1)$-set boundaries, by
\Uref{ULCD.SimplexMin}. The source's needed conclusion is LCD for all
$k$, which now has a uniform proof.
Its formula on physical page~15 is not literally correct: at $k=2$
the first printed sum is empty and $k\mathbf j=0$, leaving only
$A^2e_x$. At a different pair meeting $x$ in one point this has
coefficient one, whereas $e_x$ has coefficient zero.
Equation~\eqref{eq:odd-split} supplies the missing complementary term.
This refutes that formula as printed in F017, not the LCD conclusion
and not an unexamined assertion in the cited Fish--Fray--Mwambene paper.

\subsection{All Hamming examples, with their scope}
Write $H_k(d,q)$ for the graph on $q$-ary words of length $d$, joined
when their Hamming distance is $k$. The common distance-count formula
\Uref{ULCD.HammingCounts} and scalar test cover the concluding examples.
For odd $q$, the one-coordinate off-diagonal operators are commuting
idempotents over $\mathbb F_2$, and each distance-$k$ operator is the
corresponding elementary symmetric polynomial. Hence every
$H_k(d,q)$ is binary RLCD by \Uref{ULCD.HammingOddAlphabet}.
In particular this proves both reported $H_1(6,3)$ and $H_2(6,3)$ cases.

For $H_2(d,2)$ in characteristic three the only nonzero common-neighbour
counts can occur at distances $0,2,4$, and are
$\binom d2$, $2(d-2)$ and $6$ respectively. Adding the adjacency term
at distance two makes all entries zero when $d\equiv0\pmod3$.
Thus this entire family is RLCD, including $d=6,9$ from the source.
No unverified classification of all primes, alphabets or distances is
inferred from these examples.

For the binary cube $H_1(d,2)$ with positive even $d$, the common cube
theorem gives $A^2=0$, rank $2^{d-1}$, and
$C_A=C_A^\perp$. Its complement matrix is invertible:
$B^2=I$ in characteristic two. Consequently the printed complementary
duality assertion for $d=6,8,10$ must be corrected to self-duality of
the adjacency code. It is not an RLCD or CLCD example.

The concluding references to McLaughlin subgraphs and Hamming papers
are source context. The former is the separate L050 application, with
its own full finite proof obligations. The latter includes an
``in preparation'' item in the 2017 bibliography; we do not turn that
historical status into a present-day open-problem claim.
The source's prospective search for more RLCD and CLCD families and
for efficient decoders remains outside the proved contract except for
the uniform families explicitly established above.

\section{Attribution, finite premises and formal scope}
All graph and code results are attributed to the source and the
predecessors it names. The common paper gives complete proofs of the
projector, scalar-count, binary incidence, ternary dimension and
transitive-enumeration arguments. Its complete subset-residue certificate
and the two complete Paley programs and their records remain actual
finite premises. Optional small checks of the local identities and
counterexamples are inventoried separately. None of these programs is
silently omitted from the proof boundary because a written-text
comparison excludes software symmetrically.

The common six generic projection declarations and three generic
transitive-distance deductions are the existing scoped Lean evidence.
They do not formalize the graph families, the source-specific incidence
bridges, retained spectra and distances, or the finite computations.
No further formalization is claimed by this application. The exact
source-contract and dependency records accompany the paper; unknown
external proof, historical software and graphic costs are not treated
as zero.

\begin{thebibliography}{9}
\bibitem{source} J. D. Key and B. G. Rodrigues,
\emph{LCD codes from adjacency matrices of graphs}, author manuscript,
4 August 2017, 17 pp. (F017).
\bibitem{scheme} A. E. Brouwer and W. H. Haemers,
\emph{Association schemes}, Handbook of Combinatorics, vol. 1,
Elsevier/MIT, 1995, pp. 749--771; P. Delsarte,
\emph{An algebraic approach to the association schemes of coding theory},
Philips Research Reports, Supplement 10, 1973.
\bibitem{spectrum} M. Krebs and A. Shaheen,
\emph{On the spectra of Johnson graphs}, Electron. J. Linear Algebra
17 (2008), 154--167.
\bibitem{triples} J. D. Key, J. Moori and B. G. Rodrigues,
\emph{Binary codes from graphs on triples}, Discrete Math. 282 (2004),
171--182; \emph{Ternary codes from graphs on triples}, Discrete Math.
309 (2009), 4663--4681; \emph{Permutation decoding for binary codes
from triangular graphs}, European J. Combin. 25 (2004), 113--123.
\bibitem{incidence} W. Fish, J. D. Key and E. Mwambene,
\emph{Codes from adjacency matrices of uniform subset graphs}, cited
as submitted in F017; W. Fish, \emph{Codes from uniform subset graphs
and cycle products}, Ph.D. thesis, University of the Western Cape, 2007.
\bibitem{reflexive} W. Fish, J. D. Key and E. Mwambene,
\emph{Ternary codes from reflexive graphs on 3-sets}, Appl. Algebra
Engrg. Comm. Comput. 25 (2014), 363--382; \emph{Binary codes from
reflexive graphs on 3-sets}, Adv. Math. Commun. 9 (2015), 211--232.
\bibitem{odd} W. Fish, R. Fray and E. Mwambene,
\emph{Binary codes and partial permutation decoding sets from the odd
graphs}, Cent. Eur. J. Math. 12 (2014), 1362--1371.
\bibitem{paley} P. Seneviratne and J. Limbupasiriporn,
\emph{Permutation decoding of codes from generalized Paley graphs},
Appl. Algebra Engrg. Comm. Comput. 24 (2013), 225--236;
J. Limbupasiriporn and P. Limbupasiriporn,
\emph{Codes from adjacency matrices of generalized Paley graphs},
cited as in preparation in F017.
\bibitem{massey} J. L. Massey, \emph{Linear codes with complementary
duals}, Discrete Math. 106/107 (1992), 337--342.
\end{thebibliography}
